% !TEX program = pdflatex
\documentclass[3p,11pt]{elsarticle}
\usepackage{amsmath,amssymb}
\usepackage{array}
\usepackage{booktabs}
\usepackage[tableposition=top]{caption}
\usepackage{enumitem}
\usepackage{graphicx}
\usepackage{tikz}
\usepackage{pgfplots}
\pgfplotsset{compat=1.18}
\usepackage{url}
\urlstyle{same}
\usepackage[hidelinks]{hyperref}
\emergencystretch=3em
% first-page footer: the page number alone
\makeatletter
\def\ps@pprintTitle{\let\@oddhead\@empty\let\@evenhead\@empty
  \def\@oddfoot{\hfil\thepage\hfil}\let\@evenfoot\@oddfoot}
\makeatother

% ---- notation ---------------------------------------------------------------
\newcommand{\kms}{km\,s$^{-1}$}
\newcommand{\ms}{m\,s$^{-1}$}
\newcommand{\Msun}{M$_\odot$}
\newcommand{\Mjup}{M$_\mathrm{J}$}
\newcommand{\Mearth}{M$_\oplus$}

\begin{document}

\begin{frontmatter}

\title{Companions hidden in public HARPS velocities: four stellar companions with no published orbit}

\author{Charles C. Norton}

\begin{abstract}
HARPS-RVBank holds 212\,562 public velocities of 2\,912 stars, and Gaia DR3 measures the same
stars' motions on the sky. Searched together, they reveal companions that no catalogue holds. Every
star whose velocities vary as a companion's would is set against SIMBAD, the papers it lists, Gaia
DR3's two-body orbits and the NASA Exoplanet Archive, and the rest are fitted with one Keplerian
orbit, weighed where the orbit outlasts the velocities against the Hipparcos-Gaia proper-motion
anomaly, Gaia's acceleration and trend solutions and the range of its own epoch velocities. Four
stellar companions follow with no published orbit: a 0.41\,\Msun\ star seen nearly face on in a
465\,d orbit of the F5 dwarf HD\,7399, whose velocities alone give a brown dwarf's minimum mass; a
0.21\,\Msun\ star in a 26\,yr orbit of the K7 dwarf GJ\,4001; a companion of at least
0.30\,\Msun\ in an eccentric 344\,d orbit of the F7 dwarf HD\,77462; and a star of
$1.15 \pm 0.07$\,\Msun\ in a $27.2^{+2.0}_{-1.5}$\,yr orbit of the 2.92\,\Msun\ G8 giant
HD\,20037, with an eccentricity of $0.45 \pm 0.02$ and an inclination of $98.4 \pm 1.1^\circ$, from
the giant's HARPS velocities, Gaia DR3's velocity-trend and acceleration solutions and its
Hipparcos-Gaia proper-motion anomaly fitted together. The giant's spectra hold G-type lines where
that companion's track predicts them, at 1.7 per cent of a G0 dwarf's line depth and $5.4\sigma$
against permuted tracks.
\end{abstract}

\begin{keyword}
Techniques: radial velocities \sep Astrometry \sep Binaries: spectroscopic \sep Stars: individual: HD\,20037
\end{keyword}

\end{frontmatter}

\section{Introduction}\label{sec:intro}

Radial-velocity surveys have watched thousands of nearby stars for two decades, and their
archives hold more than their programmes published. HARPS-RVBank~\cite{Trifonov2020} recomputed
every public HARPS velocity with SERVAL~\cite{Zechmeister2018}, corrected the instrument's
nightly zero points, and attached each spectrum's activity indicators, so that a star's whole
record can be read at once. Gaia DR3 adds two independent views of the same stars: its own
spectroscopic velocities and their variability, with orbital and trend solutions for some, and
astrometry whose acceleration and proper-motion anomaly against Hipparcos~\cite{Kervella2022} bound
a companion's mass and inclination~\cite{GaiaNSS2023}.

This paper reads the two together for companions that no catalogue holds. Section~\ref{sec:data}
describes the data, Section~\ref{sec:stellar} the stellar and substellar companions and
Section~\ref{sec:hd20037} the companion of the giant HD\,20037.

\section{Data}\label{sec:data}

The velocities are HARPS-RVBank's SERVAL velocities with their nightly zero points corrected,
2\,912 stars and 212\,562 velocities, with the chromatic index, differential line width,
H$\alpha$ and Na\,D indices and the pipeline's line width and bisector of each spectrum. Each
star is matched to Gaia DR3~\cite{GaiaNSS2023} as the brightest source within 30\,arcsec of its
J2000 position whose $G$ lies within 2.5\,mag of its $V$, which follows stars of high proper
motion to Gaia's epoch, and to the Hipparcos-Gaia proper-motion anomaly of
ref.~\cite{Kervella2022} through its Hipparcos number, with that catalogue's mass of the star. The
public HARPS spectra of the stars discussed here come from the ESO Science Archive, and their
ultraviolet and X-ray fluxes from GALEX~\cite{Bianchi2017}, eROSITA's first
survey~\cite{Merloni2024}, ROSAT's 2RXS~\cite{Boller2016} and 4XMM. Every analysis is that of
orbitlight (\url{https://huggingface.co/datasets/phanerozoic/orbitlight}), whose tables give each
star's decision and each orbit's posterior draws.

\section{Stellar and substellar companions}\label{sec:stellar}

A star enters the search when its nightly velocities vary as a companion's would, with a reduced
$\chi^2$ above 25 and a scatter above 30\,\ms\ and above five times the median error on four
nights or more, and when it is of spectral type F, G, K or M and luminosity class III or fainter;
the signals of the planet search with a minimum mass above 13\,\Mjup\ or a semi-amplitude of
100\,\ms\ or more join them. Of 341 stars that vary, 273 are known: SIMBAD lists each as a binary
or as a host of planets or brown dwarfs, a paper on a binary, a companion or an orbit covering at
most 500 objects names it, Gaia DR3 holds a two-body orbit for it, or the NASA Exoplanet Archive
lists planets around it. Another 55 have fewer than ten nights. The remaining 13, with 58 stars
whose planet-search signals are that heavy, are analysed. An eccentric periodogram gives each
star's best orbit, its false-alarm level from velocities permuted within each series, and the
orbit is sampled where the velocities cover two cycles. Where they cover fewer, one Keplerian orbit
of the star about the barycentre, the companion taken dark and no heavier than the star, is fitted
to the velocities, the Hipparcos-Gaia proper-motion anomaly~\cite{Kervella2022}, Gaia DR3's
acceleration and velocity-trend solutions and, where Gaia's spectrometer saw the velocity vary,
the range of its epoch velocities less 2.5\,\kms\ as a lower bound on the orbit's range over
Gaia's window. A profile over period, the best orbit in each of 120 bands of log
period found by differential evolution, gives the periods the data allow, those within 2 of the
best in log posterior, and the posterior is sampled by an ensemble of differential-evolution moves.
A companion is accepted where a second reduction of the archived spectra, by cross-correlation,
follows the velocities with a correlation of 0.9 or more, and either the velocities cover two
cycles with a false-alarm level below $10^{-3}$ and no other period within 10 in $\chi^2$, or the
astrometry is significant (an anomaly of $5\sigma$, or a Gaia acceleration or trend solution) and
the allowed periods span less than a factor of two clear of the prior's limits. SIMBAD's references
and the papers they cite give each star's literature.

Three companions have no published orbit (Table~\ref{tab:stellar}). HD\,7399, an F5 dwarf, moves
on a 465\,d orbit at only 0.52\,\kms, a minimum mass of 0.023\,\Msun\ in the brown-dwarf desert,
but its proper-motion anomaly of $15\sigma$, Gaia's acceleration solution (Acceleration9) and its
RUWE of 1.87 move it far more on the sky than that mass could: the orbit is seen within
$4^\circ$ of face on and the companion is a 0.41\,\Msun\ star. GJ\,4001 (MCC\,791), a K7 dwarf,
has a companion of at least 0.041\,\Msun\ on an eccentric 26\,yr orbit that its six-year series
covers in part, and its anomaly of $64\sigma$ and RUWE of 9.2 put the orbit $13^\circ$ from face
on and the companion at 0.21\,\Msun, an M dwarf. HD\,77462, an F7 dwarf, has 11 HARPS nights over
472\,d, and Gaia's epoch velocities span 20.7\,\kms; together they give an orbit of 344\,d with
an eccentricity of 0.68 and a companion of at least 0.30\,\Msun. Fitted as dark, its anomaly of
$5.8\sigma$ and RUWE of 14.0 make it 1.23\,\Msun\ at $i = 161^\circ$, which a main-sequence star
of that mass cannot be: it would give nearly as much light as the primary and move the photocentre
little, so the companion is compact or the astrometric excess has a further source.

\begin{table}[t]
\centering
\caption{Companions with no published orbit: the nights of velocities, the periods the data allow
and the orbit's period, semi-amplitude and eccentricity with their 16th to 84th percentiles, the
primary's mass $m_1$~\cite{Kervella2022}, the companion's minimum mass and its mass from the
astrometry, and Gaia's RUWE and the Hipparcos-Gaia proper-motion anomaly's signal to noise.}\label{tab:stellar}
\footnotesize\setlength{\tabcolsep}{3pt}
\resizebox{\textwidth}{!}{%
\begin{tabular}{llcccccccc}
\toprule
Star & Type & Nights & $P$ (d), allowed & $K$ (\kms) & $e$ & $m_1$ (\Msun) & $m_2\sin i$ (\Msun) & $m_2$ (\Msun) & RUWE, PMa S/N\\
\midrule
HD\,7399 & F5\,V & 22 & 465 (460--473), 441--478 & 0.523 (0.516--0.533) & 0.07 (0.05--0.12) & 1.30 & 0.023 & 0.41 (0.38--0.44) & 1.87, 15.1\\
GJ\,4001 & K7\,V & 26 & 9380 (8660--10210), 7540--11350 & 0.833 (0.815--0.852) & 0.76 (0.75--0.78) & 0.63 & 0.041 & 0.21 (0.19--0.24) & 9.24, 64.3\\
HD\,77462 & F7\,V & 11 & 344.0 (343.8--344.2), 330--359 & 8.86 (8.41--9.61) & 0.68 (0.66--0.70) & 1.30 & 0.30 (0.28--0.31) & 1.23 (1.13--1.37) & 14.0, 5.8\\
\bottomrule
\end{tabular}}
\end{table}

\section{The companion of the giant HD\,20037}\label{sec:hd20037}

HD\,20037 is a G8 giant of $2.92 \pm 0.15$\,\Msun\ and 10.3\,R$_\odot$~\cite{Kervella2022} at
172\,pc. Its 13 HARPS velocities fall smoothly by 3.3\,\kms\ from 2006 to 2011 with nights a day
apart agreeing to a few metres per second. Gaia DR3 holds three more views of the same orbit. Its
27 velocities over 2014 to 2017 vary beyond their errors and are fitted by a second-degree trend
(SecondDegreeTrendSB1) rising at $0.0071 \pm 0.0036$\,\kms\ per day at 2016.0; its astrometry
carries an acceleration solution (Acceleration7) at a significance of 31.9, $-1.468 \pm 0.046$
and $-0.705 \pm 0.047$\,mas\,yr$^{-2}$ in right ascension and declination, with the proper
motion at 2016.0; and its mean motion between Hipparcos and Gaia differs from that proper motion
at $114\sigma$~\cite{Kervella2022}.

One Keplerian orbit of the giant about the barycentre, the companion dark, is fitted to all of it
at once. The HARPS velocities are read directly with a zero point and a jitter. Gaia's velocity
derivatives are those of a quadratic fitted to the model's velocities at 27 times across its
968-day velocity window; the proper motion and acceleration at 2016.0 are those of a quadratic
fitted to the photocentre over Gaia's astrometric window; and the long-term mean motion is the
photocentre's displacement from its Hipparcos position, a line fitted over Hipparcos's window,
to its Gaia position over the 24.75\,yr between them. The parameters are the period, the time of
periastron, $\sqrt{e}\cos\omega$, $\sqrt{e}\sin\omega$, the companion's mass as a fraction of the
giant's, $\cos i$, the node, the HARPS zero point, the jitter and the giant's mass with a prior of
10 per cent; Gaia's 27 velocities span 7.7\,\kms, and less 2.5\,\kms\ for their noise that range
bounds the orbit's over Gaia's window from below. A profile over period located the posterior,
which allows periods from 8\,830 to 12\,190\,d, and an ensemble of differential-evolution moves
sampled it (Table~\ref{tab:hd20037}). The companion weighs $1.15 \pm 0.07$\,\Msun\ on an orbit of
27.2\,yr seen nearly edge on; periastron passed in 2013, between the HARPS and Gaia series.

\begin{table}[t]
\centering
\caption{HD\,20037's orbit from its HARPS velocities and Gaia's and Hipparcos's astrometry
together: medians with 16th to 84th percentiles of the sampled posterior.}\label{tab:hd20037}
\small
\begin{tabular}{lc}
\toprule
Period (d) & $9\,953$ ($9\,384$--$10\,666$)\\
Eccentricity & 0.45 (0.43--0.47)\\
Giant's semi-amplitude $K_1$ (\kms) & 4.99 (4.94--5.04)\\
Inclination (deg) & 98.4 (97.4--99.6)\\
Semi-major axis (AU) & 14.4 (13.7--15.1)\\
Companion's mass (\Msun) & 1.15 (1.08--1.21)\\
\bottomrule
\end{tabular}
\end{table}

A main-sequence star of 1.15\,\Msun, an F8 or G0 dwarf, would give 3 to 4 per cent of the
system's $V$-band light. With the orbit known, the companion's lines can be sought along their own
track: in each spectrum it moves against the giant at $-(1 + m_1/m_2)$ times the giant's velocity
about the systemic one, here $+8.1$ to $+20.0$\,\kms. The 13 spectra are placed in the giant's
frame on a grid uniform in velocity, divided by their median, and cross-correlated with the line
depths of the G0 dwarf $\zeta$\,Tuc, each wavelength weighted by the spectra's scatter smoothed
over 501\,\kms, and the sky's pattern is taken out of every row in its own frame. Each row is then
fitted as a pattern common to every row plus $\alpha$ times the G0 lines' own cross-correlation at
the companion's velocity, the two by weighted least squares together, so that $\alpha$ is the depth
of the companion's lines as a fraction of a G0 dwarf's and the part of the trail the median
spectrum holds is fitted with it. Along the track $\alpha = 0.017$, at $5.4\sigma$ against 400
permutations of the track among the spectra, and the G0 spectrum laid into the spectra along the
track at 1, 2 and 4 per cent of the light raises $\alpha$ by 0.010, 0.020 and 0.040. Along the
track mirrored about the giant's velocity $\alpha$ is 0.001, and along the telluric lines' own track
$-0.004$. The spectra see the companion only between $+8$ and $+20$\,\kms, so a component at a fixed
velocity near $+12$\,\kms\ from the giant fits them as well as one on the orbit; Gaia DR3 holds no
other source within 20\,arcsec and finds no second peak in the star's images, which would leave
such a component within about 0.2\,arcsec of the giant. GALEX measures HD\,20037 at FUV
$19.76 \pm 0.11$ and NUV $14.24 \pm 0.01$, a G8 giant's colours with no hot component, and eROSITA
records a soft source 1.8\,arcsec from the star, 17 counts mostly between 0.5 and 1\,keV,
$2.0\times10^{29}$\,erg\,s$^{-1}$ in the 0.2 to 2.3\,keV band, as a giant's corona or an active
dwarf gives.

The companion is therefore most likely a main-sequence star near 1.15\,\Msun, whose lines, at
1.7 per cent of a G0 dwarf's depth against the 3 to 4 per cent of the light such a star gives, are
shallower than a slowly rotating G0 dwarf's, as rotation or a slightly earlier type makes them.
The orbit's fit takes the companion as dark; at 1.7 per cent of the light the photocentre moves
about 6 per cent less than the giant, which the fit reads as a slightly lighter companion.

\section{Conclusion}\label{sec:conclusion}

Public HARPS velocities read beside Gaia DR3's astrometry and velocities give four stellar
companions with no published orbit: a brown-dwarf-desert minimum mass that the astrometry turns
into a 0.41\,\Msun\ star seen face on around HD\,7399, an M dwarf on a 26-year orbit of GJ\,4001,
a companion of HD\,77462 heavier than 0.30\,\Msun\ whose astrometric mass no main-sequence star
could have, and a companion of $1.15 \pm 0.07$\,\Msun\ on a 27-year orbit of the giant HD\,20037,
whose lines appear in the giant's spectra along its predicted track at 1.7 per cent of a G0
dwarf's depth.

\section*{Acknowledgements}

Based on data obtained from the ESO Science Archive Facility, and on HARPS-RVBank. This work has
made use of data from the European Space Agency mission Gaia, processed by the Gaia Data
Processing and Analysis Consortium, funded by national institutions, in particular the
institutions participating in the Gaia Multilateral Agreement; of GALEX and eROSITA data; of the NASA Exoplanet Archive,
operated by the California Institute of Technology under contract with NASA; and of the SIMBAD
database and the VizieR catalogue access tool, operated at CDS, Strasbourg, France.

\begin{thebibliography}{99}\small\raggedright
\setlength{\itemsep}{1pt}

\bibitem{Trifonov2020}
T.~Trifonov et al., Public HARPS radial velocity database corrected for systematic errors,
\emph{Astron. Astrophys.} \textbf{636}, A74 (2020).

\bibitem{Zechmeister2018}
M.~Zechmeister et al., Spectrum radial velocity analyser (SERVAL). High-precision radial
velocities and two alternative spectral indicators, \emph{Astron. Astrophys.} \textbf{609}, A12
(2018).

\bibitem{Kervella2022}
P.~Kervella, F.~Arenou and F.~Th\'evenin, Stellar and substellar companions from Gaia EDR3.
Proper-motion anomaly and resolved common proper-motion pairs, \emph{Astron. Astrophys.}
\textbf{657}, A7 (2022).

\bibitem{GaiaNSS2023}
Gaia Collaboration, F.~Arenou et al., Gaia Data Release 3. Stellar multiplicity, a teaser for the
hidden treasure, \emph{Astron. Astrophys.} \textbf{674}, A34 (2023).

\bibitem{Bianchi2017}
L.~Bianchi, B.~Shiao and D.~Thilker, Revised catalog of GALEX ultraviolet sources. I. The
all-sky survey: GUVcat\_AIS, \emph{Astrophys. J. Suppl.} \textbf{230}, 24 (2017).

\bibitem{Merloni2024}
A.~Merloni et al., The SRG/eROSITA all-sky survey. First X-ray catalogues and data release of the
western Galactic hemisphere, \emph{Astron. Astrophys.} \textbf{682}, A34 (2024).

\bibitem{Boller2016}
Th.~Boller et al., Second ROSAT all-sky survey (2RXS) source catalogue, \emph{Astron. Astrophys.}
\textbf{588}, A103 (2016).

\end{thebibliography}

\end{document}
